%% User macros

%% TeXtured logo
\newcommand*{\TeXtured}{\texorpdfstring{\textsf{\TeX{}tured}}{TeXtured}}

\usepackage{hologo} % Provides *TeX logos (for example `\hologo{pdfLaTeX}`)

%%% Game terms: one consistent style per kind (readable on a phone and in greyscale print)
%%   aspect: bold italic | skill: sans semibold | trick: italic in TrickColor
%%   action: small caps  | peripetie: sans bold
%% skills: Source Sans Semibold scaled to the x-height of the text, so a skill name does not look bigger.
%%   The font files are named differently across TeX Live versions: "Source Sans Pro" (up to 2025)
%%   became "Source Sans 3" (2026, used by the CI), so both names are tried; if neither is found,
%%   the bold sans of the headings is used instead of failing the build.
\IfFontExistsTF{SourceSansPro-Semibold.otf}{%
    \newfontfamily\dovfont{SourceSansPro-Semibold.otf}[Scale=MatchLowercase]%
}{\IfFontExistsTF{SourceSans3-Semibold.otf}{%
    \newfontfamily\dovfont{SourceSans3-Semibold.otf}[Scale=MatchLowercase]%
}{%
    \NewDocumentCommand{\dovfont}{}{\sffamily\bfseries}%
}}
\newcommand{\asp}[1]{\asteriontermzapis{asp}{#1}\textbf{\textit{#1}}}       % aspect
\newcommand{\dov}[1]{{\dovfont #1}}              % skill
\newcommand{\trk}[1]{\asteriontermzapis{trk}{#1}\textcolor{TrickColor}{\textit{#1}}} % trick (no underline: on screen it looks like a link)

%%% Lists of all aspects and tricks in the appendix (`chapters/dodatky/rejstrik.tex`, Lua side in
%%   `preamble/user/rejstrik.lua`): every \asp/\trk records its text and page in the .aux file,
%%   the next run prints them sorted, each term once with all its pages
\directlua{dofile("preamble/user/rejstrik.lua")}
\makeatletter
\NewDocumentCommand{\asteriontermzapis}{mm}{%
    \protected@write\@auxout{\let\thepage\relax}{\string\asteriontermin{#1}{\unexpanded{#2}}{\thepage}}%
}
\makeatother
\NewDocumentCommand{\asteriontermin}{mmm}{%
    \directlua{asterion_rejstrik.add("\luaescapestring{#1}", "\luaescapestring{\detokenize{#2}}", "\luaescapestring{#3}")}%
}
\pdfstringdefDisableCommands{\def\asteriontermzapis#1#2{}}% no recording inside PDF bookmarks
%% `\vypisrejstrik{asp}` / `\vypisrejstrik{trk}`: the list (two columns on A4, one on the phone)
\NewDocumentCommand{\vypisrejstrik}{m}{%
    \begin{multicols}{2}\raggedright\setlength{\parskip}{0pt plus 1pt}% compact: one line per entry
    \directlua{asterion_rejstrik.print("\luaescapestring{#1}")}%
    \end{multicols}%
}
\NewDocumentCommand{\asterionrejstrikpismeno}{m}{%
    \par\addvspace{0.6\baselineskip}%
    {\sffamily\bfseries\color{FateBlue}#1}\par\nobreak
}
\NewDocumentCommand{\asterionrejstrikpolozka}{mmm}{%
    \par\hangindent=1em\hangafter=1\noindent
    \ifstrequal{#1}{asp}{\textbf{\textit{#2}}}{\textcolor{TrickColor}{\textit{#2}}}%
    \nobreak\enspace{\footnotesize\color{black!60}#3}\par
}
\NewDocumentCommand{\asterionrejstrikprazdny}{}{\emph{(Seznam se doplní při dalším překladu.)}\par}
\newcommand{\akc}[1]{\textsc{#1}}                % action
\newcommand{\per}[1]{\textbf{\sffamily #1}}      % peripetie

%%% Image credits: `\zdrojobrazku{<url>}` under an illustration prints a small "zdroj" link
%%   (instead of a numbered caption) and records the page and the site for the list of image
%%   sources at the end of the book (`\zdrojeobrazku`, called in `main.tex`)
\NewDocumentCommand{\zdrojobrazku}{m}{%
    \par\nobreak\vspace{2pt}%
    {\centering\footnotesize\href{#1}{zdroj}\par}%
    \addtocontents{zdroje}{\protect\zdrojobrazkuradek{\thepage}{\detokenize{#1}}}%
}
%% one line of the list: page, then the site's domain as a link to the source
\NewDocumentCommand{\zdrojobrazkuradek}{mm}{%
    \par\noindent\makebox[3.5em][l]{s.~#1}%
    \href{#2}{\directlua{
        local p = string.char(37) % Lua's pattern escape (a literal percent sign would be a TeX comment)
        local url = "\luaescapestring{#2}"
        local host = string.match(url, "^" .. p .. "a+://([^/]+)") or url
        tex.sprint(-2, (host:gsub("^www" .. p .. ".", "")))
    }}\par
}
\makeatletter
\NewDocumentCommand{\zdrojeobrazku}{}{%
    \chapter*{Zdroje obrázků}%
    \addcontentsline{toc}{chapter}{Zdroje obrázků}%
    {\hypersetup{linkcolor=FateBlue}\@starttoc{zdroje}}%
}
\makeatother

\ExplSyntaxOn

% Define \policko{<number>}{<count>}
% a line may break between the boxes, so long rows wrap instead of overflowing
\NewDocumentCommand \policko { m m }
  {
    \int_step_inline:nn { #2 }
      {
        \int_compare:nNnF { ##1 } = { 1 }
          { \penalty 0 \hskip 0.25em plus 0.25em \scan_stop: }
        \fbox{#1}
      }
  }

\ExplSyntaxOff

